%=====================================================================
\chapter{Academic Integrity, Plagiarism and the Similarity Test}\label{ch:integrity}
%=====================================================================
\locator{5}{Part~5 --- Rigour, Reproducibility and Ethics}

\begin{outcomes}
By the end of this chapter you will be able to:
\begin{tight}
  \item \textbf{State} the numeric similarity thresholds that bind your thesis.
  \item \textbf{Interpret} an originality report rather than reading its
    percentage.
  \item \textbf{Paraphrase} and cite so that inadvertent plagiarism cannot
    occur.
  \item \textbf{Apply} authorship criteria, and recognise gift and ghost
    authorship.
  \item \textbf{Distinguish} text recycling, salami slicing, fabrication and
    falsification.
\end{tight}
\end{outcomes}

\begin{prereq}
\chref{appraisal} for the note-taking discipline that prevents most accidental
plagiarism, and \chref{bibliometrics} for citation management.
\chref{publication} for authorship in the context of \S15.
\end{prereq}

\begin{keyterms}
plagiarism $\bullet$ similarity index $\bullet$ originality report $\bullet$
text recycling $\bullet$ self-plagiarism $\bullet$ salami slicing $\bullet$
fabrication $\bullet$ falsification $\bullet$ authorship criteria $\bullet$
gift and ghost authorship $\bullet$ CRediT
\end{keyterms}

\section{The Rules, Quoted}

These are numeric and non-negotiable, and they are the same in both degree
regulations. Read them once carefully; the rest of this chapter is about
meeting them.

\begin{regbox}[Similarity --- PhD \S32(III), MPhil/MS \S49(3)]
``Plagiarism is not allowed at all. The similarity index should be less than or
equal to \textbf{19\%} which is supposed to be a connection with the existing
knowledge and 80\% of research is contributed by the author. Further, less than
\textbf{5\% from a single source} is allowed. The Similarity index should be
considered very seriously in the section of findings and conclusion of the
document. The similarity index for that section should not be more than
\textbf{9\%}.''

\medskip
And \S32(VI): ``A similarity test, in accordance with the HEC's Anti-Plagiarism
Policy, must be conducted on the dissertation before its submission to the
external experts by the student.''
\end{regbox}

\begin{regbox}[What counts as plagiarism --- PhD \S32(I)(A)]
``(i) Stealing and presenting the ideas or words of others as one's own;
(ii) Using another person's production, without citing and crediting the
source; (iii) Committing the literary theft; (iv) Presenting as a new and
original idea or product derived from an existing scholarly source;
(v) Turning in someone else's work as one's own; (vi) Copying words or ideas
from someone else without giving credit; (vii) Failing to put a quote or
quotation marks, when copying the exact language from a source; (viii) Giving
incorrect information about the source of a quotation; (ix) Changing words but
copying the sentence structure of a source without giving credit; (x) Copying a
bulk of words or ideas from other references and including them in your work,
whether you give credit or not.''
\end{regbox}

\begin{alertbox}[Items (ix) and (x) catch people who thought they were safe]
\textbf{(ix) Changing words but copying the sentence structure.} This is what
most students think paraphrasing is. Substituting synonyms into someone else's
sentence, with a citation attached, is still plagiarism under this clause. The
sentence must be yours, built from your understanding.

\medskip
\textbf{(x) Copying a bulk of words or ideas \emph{whether you give credit or
not}.} A correctly cited page-long quotation is still a problem: the work must
be substantially your own writing. \S32(III)'s ``80\% of research is
contributed by the author'' is the same idea expressed as a proportion.
\end{alertbox}

\section{Reading an Originality Report}

The similarity index is not a plagiarism score. The regulations themselves say
so, and knowing how to read the report is the practical skill.

\begin{regbox}[Interpretation --- PhD \S29(V)]
``The similarity index of the originality report showing matches of submitted
work with internet content does not mean that the document is plagiarized. The
similarity index is based on the percentage of matched text out of the total
number of words in a document\ldots{} Common phrases and proper nouns also
appear as similarities in the report; therefore, every instructor or faculty
member should ignore the matches returned from them\ldots{} The graphs, tables,
formulae, and other pictorial materials are not matched through the service;
therefore, they will only offer similarities with text.''

\medskip
And \S29(IV)(iv): ``References/bibliographies and tables of contents must be
removed from the submitted documents. If these are included, the similarity
index of the document will increase.''

\medskip
\S29(IV)(viii): ``If the report has a similarity index of $\le$19\%, then the
benefit of the doubt may be given to the author; however, if any single source
has a similarity index $\ge$5\% without citations then it needs to be revised.''
\end{regbox}

\begin{center}
\small
\begin{tabular}{@{}L{4.4cm}L{9.9cm}@{}}
\toprule
\textbf{What you see} & \textbf{What it means} \\
\midrule
18\% overall, spread thinly across hundreds of sources
  & Normal. Method names, standard phrases, technical terms. Check a sample and
    move on \\
8\% overall, but 7\% from one source
  & \textbf{A problem}, despite the low total. Breaches the single-source rule
    and is what genuine copying looks like \\
High match in the references section
  & An artefact. References should have been excluded before submission \\
High match in Methods
  & Often benign --- standard procedures are described in standard language ---
    but check that it is paraphrase and not lifted text \\
High match in Findings or Conclusions
  & \textbf{Serious.} These must be your own words about your own results, and
    the 9\% cap applies specifically here \\
Match to your own earlier paper
  & Text recycling. Discussed below; \S29(V)(iv) permits ignoring it if the
    material is cited \\
\bottomrule
\end{tabular}
\end{center}

\begin{conceptbox}[The index you should aim for is not zero]
A thesis with 2\% similarity is not obviously excellent; it may indicate that
standard terminology has been avoided in ways that make the writing worse, or
that the engagement with prior work is thin.

\medskip
The target is a report in which \emph{every match is explicable}: this is a
quoted and cited passage; this is a standard method name; this is the
university's own thesis template. Go through the report match by match before
submission and be able to explain each. That is what the supervisor's verdict
under \S29(V)(ix) requires, and preparing for it is straightforward.
\end{conceptbox}

\section{Paraphrasing Properly}

\begin{center}
\small
\begin{tabular}{@{}L{4.6cm}L{9.7cm}@{}}
\toprule
\textbf{Original} & ``Systematic reviews aggregate evidence across primary
                    studies using an explicit, reproducible protocol that
                    reduces the influence of reviewer bias.'' \\
\midrule
\textbf{Not acceptable} --- (ix)
                  & ``Systematic reviews combine evidence from primary studies
                    using a clear, repeatable protocol that lessens the effect
                    of reviewer bias [12].'' \emph{Same structure, synonyms
                    substituted.} \\
\midrule
\textbf{Acceptable}
                  & ``The value of the protocol in a systematic review is that
                    it is written before the evidence is seen; that ordering,
                    rather than any property of the aggregation itself, is what
                    limits the reviewer's discretion [12].'' \emph{Different
                    sentence, and it says something the original did not ---
                    which is the sign that you understood it.} \\
\midrule
\textbf{Also acceptable}
                  & Quote it, in quotation marks, with a page number. Sparingly \\
\bottomrule
\end{tabular}
\end{center}

\begin{alertbox}[The method that makes this reliable]
Read the passage. Close the source. Wait a moment. Write what it meant, in your
own words, without looking. Then reopen and check you have not distorted it.

\medskip
You cannot paraphrase with the original sentence in front of you --- its
structure will survive, and clause (ix) catches exactly that. Combined with the
note-taking discipline from \chref{appraisal}, where quotations are marked as
quotations at the moment they are written, this eliminates essentially all
inadvertent plagiarism.
\end{alertbox}

\section{Recycling and Slicing Your Own Work}

\begin{center}
\small
\begin{tabular}{@{}L{3.2cm}L{5.2cm}L{5.7cm}@{}}
\toprule
 & \textbf{What it is} & \textbf{Position} \\
\midrule
Text recycling
  & Reusing your own text in a new work
  & Methods sections describing the same protocol are widely tolerated; results
    and discussion are not. Cite the earlier work. \S29(V)(iv) permits
    ignoring self-matches ``if it is the same work'' \\
Duplicate publication
  & The same study published twice as new
  & Never acceptable. It corrupts meta-analyses by double-counting \\
Salami slicing
  & One study divided into several minimal papers
  & Discouraged. The test: does each paper answer a distinct question and stand
    alone? If the reader needs all of them to understand any, it was one paper \\
Extended conference version
  & A journal paper extending a conference paper
  & Standard and legitimate, and often necessary under \S15
    (\chref{publication}). Disclose the earlier version, cite it, and ensure
    the extension is substantial --- typically 30\% or more new material \\
\bottomrule
\end{tabular}
\end{center}

\begin{conceptbox}[Salami slicing has a specific incentive here]
\reg{PhD Regs 2024}{\S15} requires either one W-category paper or two
X-category papers. A scholar with one good study and access only to X-category
venues faces a direct incentive to split it in two.

\medskip
Resist it, and plan around it instead: design a programme where a systematic
review (\chref{slr}) and the empirical study it motivates are genuinely
separate contributions. That satisfies the regulation with two papers that each
stand alone, which is what the regulation intends.
\end{conceptbox}

\section{Authorship}

\begin{center}
\small
\begin{tabular}{@{}L{3.6cm}L{10.7cm}@{}}
\toprule
\textbf{Authorship} (ICMJE~\cite{icmje2024recommendations}) requires all four:
  & substantial contribution to conception, design, acquisition, analysis or
    interpretation; drafting or critically revising; final approval; and
    accountability for the work's integrity \\
\textbf{Gift authorship}
  & Adding someone who did not contribute --- a head of department, a
    colleague. Common, and a breach \\
\textbf{Ghost authorship}
  & Omitting someone who did contribute. Most often a junior scholar or a
    technician \\
\textbf{CRediT}~\cite{brand2015beyond}
  & A 14-role taxonomy --- conceptualisation, methodology, software,
    validation, formal analysis, investigation, writing, supervision, and
    others. Increasingly required, and it makes contribution visible rather
    than inferred from order \\
\bottomrule
\end{tabular}
\end{center}

\begin{alertbox}[Agree authorship in writing, at the start]
\reg{PhD Regs 2024}{\S15(ii)} requires that ``the PhD researcher shall be the
first author'' of the papers counting towards the degree. A paper on which your
supervisor is first author does not count for you, however much you
contributed.

\medskip
Have the conversation before the work starts, record the agreement in an email,
and revisit it if contributions change. This is not distrust; it is the
standard professional practice, and the awkwardness of raising it early is
trivial compared with the awkwardness of raising it at submission.
\end{alertbox}

\section{Fabrication and Falsification}

\begin{center}
\small
\begin{tabular}{@{}L{3.0cm}L{11.3cm}@{}}
\toprule
Fabrication   & Inventing data or results that were never obtained \\
Falsification & Manipulating data, images, equipment or processes so that the
                research is not accurately represented --- including selective
                omission of inconvenient results \\
Image manipulation & Adjusting a figure so it shows something the data does
                not. Adjusting brightness across a whole image is acceptable;
                doing it to a region is not \\
\bottomrule
\end{tabular}
\end{center}

These are career-ending and unambiguous. The boundary worth attending to is
that \textbf{selective omission is falsification}: dropping the experimental
condition that did not work, or the three datasets where the method lost, is
not tidying. It is misrepresenting the research, and \chref{prereg} exists
partly to make it structurally difficult.

\begin{pitfall}[Where honest scholars come closest to the line]
\begin{tight}
  \item Reporting the ablation that worked and omitting the two that did not.
  \item ``Cleaning'' outliers whose removal happens to help (\chref{prereg}).
  \item Rounding a p-value of 0.054 to ``0.05'' or describing it as
    ``approaching significance''.
  \item Reusing a figure from a previous paper without saying so.
  \item Citing a paper for a claim it does not make, because a secondary source
    said it did (\chref{bibliometrics}).
\end{tight}
\end{pitfall}

\section{Consequences}

\reg{PhD Regs 2024}{\S32(V)} is unambiguous: ``Under no circumstances shall a
dissertation based on plagiarized research be acceptable. It is the primary
responsibility of both PhD researchers and their supervisors to prevent
plagiarism.'' Cases are handled under the HEC Anti-Plagiarism
Policy~\cite{hec2023plagiarism} and COPE guidance~\cite{cope2019corepractices}.

Note the shared responsibility: the supervisor is named alongside the scholar.
That is a reason to involve them early --- run the similarity check well before
submission, together, when there is still time to fix what it finds.

\begin{traditions}{integrity}
\tradEMP{Report every condition run and every analysis performed. Selective
reporting is the field's characteristic integrity failure, and it is invisible
without pre-registration.}
\tradDSR{Do not present an artefact's capabilities beyond what was built and
evaluated. Screenshots of features that do not work are falsification.}
\tradQUAL{Quotations must be verbatim, attributed, and not composited from
several participants. The audit trail is the defence (\chref{qda}).}
\tradFORM{A proof with a gap presented as complete is falsification. If a step
is unproved, say so.}
\end{traditions}

\begin{lab}[Run the check early]
\begin{tightnum}
  \item Take a draft chapter and run it through your institution's similarity
    service, with references and table of contents removed as \S29(IV)(iv)
    requires.
  \item Go through the report match by match. Classify each: quoted and cited /
    standard terminology / template / \textbf{problem}.
  \item For every problem, rewrite using the close-the-source method.
  \item Check the single-source column. Any source above 5\% without citation
    must be fixed regardless of the total.
  \item Write the authorship agreement for your intended first paper, with
    CRediT roles, and send it to your co-authors.
\end{tightnum}
\end{lab}

\begin{checkpoint}
\begin{tightnum}
  \item A thesis reports 12\% overall similarity, with 9\% from one source.
    Does it pass? Which clause decides?
  \item Why is substituting synonyms into someone else's sentence plagiarism
    even with a citation?
  \item Distinguish an acceptable extended conference version from duplicate
    publication.
  \item Your supervisor asks to be first author on a paper you wrote. What does
    \S15 imply, and how do you raise it?
\end{tightnum}
\end{checkpoint}

\begin{chaptersummary}
\begin{tight}
  \item Three numbers bind you: 19\% overall, 5\% from any single source, 9\% in
    findings and conclusions. Remove references before checking.
  \item The index is not a verdict. Aim for a report in which every match is
    explicable, and be ready to explain each.
  \item Changing words while keeping sentence structure is plagiarism under
    \S32(I)(A)(ix). Close the source, wait, write from understanding.
  \item Mark quotations as quotations when you take the note. This prevents
    nearly all inadvertent plagiarism.
  \item Text recycling in methods is tolerated; duplicate publication never is;
    salami slicing is discouraged and \S15 creates an incentive for it that you
    should plan around.
  \item Authorship requires all four ICMJE criteria. \S15 requires you to be
    first author. Agree it in writing, early.
  \item Selective omission of results is falsification, not tidying.
  \item The supervisor shares responsibility, so run the check early and
    together.
\end{tight}
\end{chaptersummary}

\begin{reviewq}
  \item Similarity thresholds are numeric and therefore auditable, but they
    measure text overlap rather than intellectual theft. Construct a thesis
    that passes all three thresholds and is nevertheless plagiarised, and one
    that fails a threshold and is not.
  \item Text recycling in methods sections is widely tolerated. Is the
    tolerance principled, or merely convenient? What rule would you propose?
  \item \S15's publication requirement creates an incentive to salami slice.
    Design a modification to the requirement that preserves its intent while
    removing the incentive.
  \item Supervisors share responsibility for plagiarism under \S32(V) yet
    typically see only late drafts. What supervisory practice would make the
    shared responsibility meaningful rather than nominal?
\end{reviewq}
